\section{Step 45: Publish loaded bytes and call the indexer}
Return the saved status and an initial count of zero. If status is success,
store the positive accumulated count at result-context offset zero and return it
in \texttt{\$v1}. Retain the context suffix and every context byte on failure.
The supplied epilogue restores saved caller state.

Fill \texttt{io\_publish}. The new caller passes the exact loaded window to
\texttt{video\_index}; it never indexes a failed or unfinished load. For an
invalid late command or checksum, loading can publish a byte count while indexing
rejects the format and retains its own directory/context outputs. This separates
file-transfer evidence from validation evidence. For replacement loading, use
separate candidate storage if an earlier accepted animation must remain usable.

\paragraph{Observe} Prepare \texttt{s45-1} and stop before
\texttt{io\_published}. At \texttt{io\_published}, \texttt{\$v0} is zero,
\texttt{\$v1} and the native count field are 24\,795, and the suffix is
unchanged. Continuing the default caller validates the file and prints
\texttt{Video file indexed.}

\Needspace{7\baselineskip}
\paragraph{Check} When may the caller invoke the indexer?
\begin{itemize}
\item[A.] After any nonnegative read.
\item[B.] Before the EOF probe.
\item[C.] After successful completed loading.
\end{itemize}
